\documentclass[11pt,reqno]{amsart}
\usepackage[T1]{fontenc}
\usepackage{lmodern}
\usepackage[margin=1in]{geometry}
\usepackage{amsmath,amssymb,mathtools,microtype,booktabs}
\usepackage[hidelinks]{hyperref}
\hypersetup{pdftitle={A short classification of complex rank-two Grothendieck rings},pdfauthor={}}
\newtheorem{theorem}{Theorem}[section]
\newtheorem{lemma}[theorem]{Lemma}
\newtheorem{corollary}[theorem]{Corollary}
\theoremstyle{remark}
\newtheorem{remark}[theorem]{Remark}
\newcommand{\C}{\mathbb C}
\newcommand{\Q}{\mathbb Q}
\newcommand{\Z}{\mathbb Z}
\newcommand{\one}{\mathbf1}
\newcommand{\cC}{\mathcal C}
\newcommand{\cS}{\mathcal S}
\newcommand{\cZ}{\mathcal Z}
\newcommand{\FPdim}{\operatorname{FPdim}}
\newcommand{\Irr}{\operatorname{Irr}}
\title{A short classification of complex rank-two Grothendieck rings}
\subjclass[2020]{18M20, 18M15}
\keywords{fusion category, rank two, Grothendieck ring, sphericalization, Drinfeld center}
\begin{document}
\begin{abstract}
We give a short proof that a complex rank-two fusion category has Grothendieck ring
$X^2=1+nX$ with $n=0$ or $1$. A norm argument in its sphericalization first
establishes sphericality without assuming pivotality. The center's Gauss sum then
excludes $n\geq2$ by a two-term triangle inequality. The argument uses standard
structural results for fusion categories and does not classify tensor structures.
\end{abstract}
\maketitle
\par\medskip\noindent\textbf{Research notice.}
This is AI-generated research. The results, proofs, code and data have not yet
been independently verified by human experts or formalized in Lean. Treat them
as provisional; automated checks establish only their stated computational scope.
\par\medskip

\section{Introduction}

Ostrik classified the complex fusion categories of rank two, including their tensor
structures \cite{OstrikRank2}. We give a compact proof at the level of Grothendieck
rings. The new proof separates the existence of a spherical structure from the
fusion-rule exclusion and does not assume a braiding on the original category.
For a fusion category with simples $\one,X$, rigidity gives $X^*=X$ and
\begin{equation}\label{eq:rule}
 X^2=\one+nX,\qquad n\in\Z_{\geq0}.
\end{equation}
Write $K_n$ for this based ring.

\begin{theorem}\label{thm:main}
For a fusion category $\cC$ over $\C$ with exactly two simple objects,
$K_0(\cC)$ is isomorphic to $K_0$ or $K_1$. Every such category admits a
spherical structure. Both rings occur, respectively in a pointed category and
in a Fibonacci category.
\end{theorem}

The statement concerns fusion rings, rather than the monoidal-equivalence
classification of their realizations. In particular, the two categories with
ring $K_0$ and the two with ring $K_1$ in \cite{OstrikRank2} need not be
identified to prove Theorem~\ref{thm:main}.

\section{A spherical structure from a norm}

We recall the general results used below. A fusion category over $\C$ admits a
\emph{sphericalization} $\pi:\cS\to\cC$, where $\cS$ is spherical. Every simple
$V$ has two simple lifts $V_+,V_-$; the lifts of $\one$ are $\one$ and a central
invertible object $\varepsilon$ of dimension $-1$, and
$\varepsilon\otimes V_+=V_-$ \cite[Remark~3.1]{ENO}. If $V$ is self-dual, so
are both lifts \cite[Section~5.1]{OstrikRank3}. We use the general
\emph{formal-codegree divisibility theorem}: if $f$ is a formal codegree of a
complex fusion category $\mathcal A$, then $\dim(\mathcal A)/f$ is an algebraic
integer \cite[Corollary~2.14]{OstrikRank3}. The global dimension here is the
intrinsic global dimension and does not require a pivotal structure.

\begin{lemma}\label{lem:lift}
Suppose $n\geq1$ in \eqref{eq:rule}. If $x$ is either lift of $X$ in the
sphericalization, its signed fusion rule modulo $\varepsilon=-1$ is
\begin{equation}\label{eq:signed}
 x^2=1+\widetilde n x,
 \qquad \widetilde n\in\Z,\quad
 |\widetilde n|\leq n,\quad\widetilde n\equiv n\pmod2.
\end{equation}
Moreover $|\widetilde n|=n$.
\end{lemma}

\begin{proof}
As $x$ is self-dual, its square contains $\one$ exactly once and
$\varepsilon$ zero times. The other constituents have the form $r x+s\varepsilon x$
with $r,s\geq0$ and $r+s=n$. Reducing modulo $\varepsilon=-1$ gives
$\widetilde n=r-s$, proving the first assertions.

Let $d=\dim(x)$. The spherical dimension character kills $\varepsilon+1$,
so $d^2=1+\widetilde n d$. The global dimension of $\cC$ equals that of one
choice of lift squared, namely
\begin{equation}\label{eq:delta}
 \delta=1+d^2;
\end{equation}
the global dimension of $\cS$ is $2\delta$. Put
$R=\Q[t]/(t^2-nt-1)$. Its two formal codegrees are $f_{\pm}=1+u_{\pm}^2$,
where $u_\pm=(n\pm\sqrt{n^2+4})/2$. Consequently
\begin{equation}\label{eq:norms}
 N_{\Q(\sqrt{n^2+4})/\Q}(f_+)=n^2+4.
\end{equation}
The number $n^2+4$ is nonsquare for $n\geq1$, so $R$ is a quadratic field.
The two lifts of each simple double each codegree of $R$ in $K_0(\cS)$.
Thus $2f_+$ is a formal codegree of $\cS$, and general codegree
divisibility says that $\delta/f_+$ is an algebraic integer.

To take its norm in a single quadratic field, we recall a second general
formal-codegree fact: the codegrees of a spherical category belong to the
field generated by its simple dimensions \cite[Corollary~2.15]{OstrikRank3}.
This field is $L=\Q(d)$, of degree at most two, and it contains both roots
$2f_\pm$ (the conjugate codegrees of $\cS$). Hence
$L=\Q(\sqrt{n^2+4})$; it is also the field of $\widetilde n$'s quadratic
character, so $\widetilde n\ne0$. The conjugate of $\delta$ is
$1+{d'}^2$ with $d'$ the other root of
$t^2-\widetilde n t-1$. Taking norms in $L$ yields
\[
 N_{L/\Q}(\delta/f_+)
 =\frac{\widetilde n^2+4}{n^2+4}\in\Z_{>0}.
\]
The ratio is positive and at most one by $|\widetilde n|\leq n$.
Therefore it equals one and $|\widetilde n|=n$.
\end{proof}

\begin{corollary}\label{cor:spherical}
Every complex rank-two fusion category admits a spherical structure.
\end{corollary}
\begin{proof}
For $n=0$ the category is pointed and hence pseudo-unitary; its canonical
pivotal structure is spherical \cite[Propositions~8.23--8.24]{ENO}.
For $n\geq1$, changing $x$ to $\varepsilon x$ changes
$\widetilde n$ to $-\widetilde n$. By Lemma~\ref{lem:lift} we may choose
$\widetilde n=n$. Since $r+s=n$ and $r-s=n$, we have $s=0$:
$x^2=\one\oplus n x$ already inside $\cS$.
The full replete subcategory with simples $\one,x$ is closed under tensor
products and duals. The forgetful tensor functor $\pi$ restricts to a tensor
equivalence from it to $\cC$. Restricting and transporting the spherical
structure of $\cS$ proves the claim.
\end{proof}

\section{The center's two known summands}

Let $\cC$ now have the spherical structure of Corollary~\ref{cor:spherical},
write $d=\dim(X)$, and set $D=\dim(\cC)=1+d^2$. We recall two standard
properties of its modular Drinfeld center $\cZ(\cC)$:
\begin{equation}\label{eq:centre}
 \sum_{B\in\Irr\cZ(\cC)}\dim(B)^2=D^2,
 \qquad
 \sum_{B\in\Irr\cZ(\cC)}\theta_B\dim(B)^2=D.
\end{equation}
The first follows from $\dim\cZ(\cC)=D^2$ and the second is the center
Gauss-sum identity \cite[Theorem~1.2]{Muger}. For completeness, the
second identity follows by writing
$\dim(B)=\sum_V[F(B):V]\dim(V)$, using adjunction, and summing
$\dim(V)\operatorname{Tr}(\theta_{I(V)})$ over the simple objects $V$:
the twist trace is $D$ for $V=\one$ and zero otherwise
\cite[Theorem~2.5]{OstrikRank3}. Here $\theta_B$ is the
canonical twist; its values are roots of unity by the Vafa theorem.

Let $F:\cZ(\cC)\to\cC$ be forgetful and $I$ its right adjoint. The
standard formula $FI(\one)=\bigoplus_V V\otimes V^*$ gives
$FI(\one)=2\one\oplus nX$. By adjunction,
$\dim_\C\operatorname{End}_{\cZ(\cC)}I(\one)=2$.
The simple unit occurs once in $I(\one)$, so semisimplicity implies
\begin{equation}\label{eq:twosummands}
 I(\one)=\one\oplus A,
 \qquad F(A)=\one\oplus nX,
 \qquad\dim(A)=1+nd=d^2.
\end{equation}
The twist of each summand of $I(\one)$ is one by the same
trivial-twist theorem \cite[Theorem~2.5]{OstrikRank3}; equivalently,
the canonical connected commutative separable algebra $I(\one)$
has trivial twist.

\begin{proof}[Proof of Theorem~\ref{thm:main}]
The sphericality assertion is Corollary~\ref{cor:spherical}. Assume $n\geq1$.
The two roots of $d^2=1+nd$ are
$d_\pm=(n\pm\sqrt{n^2+4})/2$. Choose an automorphism $\sigma$ of
$\overline\Q$ sending the actual categorical dimension $d$ to $d_+>0$.
For any simple central object $B$, its dimension is the dimension of
$F(B)=a_B\one\oplus b_B X$ with $a_B,b_B\in\Z_{\geq0}$, and hence
$\sigma(\dim(B))=a_B+b_Bd_+>0$. Applying $\sigma$ to
\eqref{eq:centre} preserves the equalities and sends each twist to
another root of unity; no Galois descent of associators is involved.

Set $t=d_+^2$ and $D_+=1+t$. The two objects in
\eqref{eq:twosummands} contribute $1+t^2$ to both identities in
\eqref{eq:centre}. Thus the other simple central objects satisfy
\[
 \sum_{B\ne\one,A}\sigma(\dim(B))^2=D_+^2-1-t^2=2t,
 \quad
 \sum_{B\ne\one,A}\sigma(\theta_B)\sigma(\dim(B))^2
 =D_+-1-t^2=t(1-t).
\]
All coefficients in the first sum are nonnegative real numbers and
$|\sigma(\theta_B)|=1$. The triangle inequality gives
$t|1-t|\leq2t$, so $t\leq3$. For $n\geq2$,
$d_+>n\geq2$ and therefore $t>4$, a contradiction. Hence $n\leq1$.
Finally $K_0$ is realized by $\mathrm{Vec}_{\Z/2}$ and $K_1$ by the
Fibonacci fusion category, for instance the even part of the
$SU(2)_3$ fusion category \cite{OstrikRank2}. Both are spherical.
\end{proof}

\begin{remark}[Positive characteristic]\label{rem:positive}
The argument above is specific to characteristic zero. Etingof, Ostrik and
Venkatesh \cite[Propositions~4.6--4.7]{EOV} prove that over an algebraically
closed field of positive characteristic no \emph{braided} fusion category
realizes $K_n$ for $n\geq2$, and explicitly leave the general fusion-category
case open. In particular their result cannot be used to assert a full
positive-characteristic classification. A category of zero global dimension
in characteristic $p$ need not have a semisimple modular center, so the proof
of \eqref{eq:centre} does not carry over directly. We know of no
later general resolution of the nonbraided rank-two problem in positive
characteristic.
\end{remark}

\paragraph{AI disclosure.}
This manuscript, including its arguments and supporting calculations, was generated with AI assistance. Independent human verification and Lean formalization are pending.

\begin{thebibliography}{9}
\bibitem{ENO}
P.~Etingof, D.~Nikshych and V.~Ostrik,
\emph{On fusion categories}, Ann. of Math. (2) \textbf{162} (2005),
581--642; \href{https://arxiv.org/abs/math/0203060}{arXiv:math/0203060}.
\bibitem{EOV}
P.~Etingof, V.~Ostrik and S.~Venkatesh,
\emph{Computations in symmetric fusion categories in characteristic $p$},
\href{https://arxiv.org/abs/1512.02309}{arXiv:1512.02309}.
\bibitem{Muger}
M.~M\"uger, \emph{From subfactors to categories and topology II:
The quantum double of tensor categories and subfactors}, J. Pure Appl.
Algebra \textbf{180} (2003), 159--219;
\href{https://arxiv.org/abs/math/0111205}{arXiv:math/0111205}.
\bibitem{OstrikRank2}
V.~Ostrik, \emph{Fusion categories of rank 2}, Math. Res. Lett.
\textbf{10} (2003), 177--183;
\href{https://arxiv.org/abs/math/0203255}{arXiv:math/0203255}.
\bibitem{OstrikRank3}
V.~Ostrik, \emph{Pivotal fusion categories of rank 3}, with an appendix
jointly with D.~Nikshych, Mosc. Math. J. \textbf{15} (2015), 373--396;
\href{https://arxiv.org/abs/1309.4822}{arXiv:1309.4822}.
\end{thebibliography}
\end{document}
